Remaining useful life (RUL) prediction must distinguish degradation from measurement failures across operating conditions. GRACE-RUL combines local CNN--GRU features, residual BiLSTM sensor recovery, Transformer context, an independent finite-fault detector and a validation-audited RUL path. Format 10 adds an observed engine-age channel and supervises continuity on the deployed repair path; it disables the bounded corrector after a null component ablation and replaces the hard repair threshold with probability-weighted soft gating. One shared model pools all four NASA C-MAPSS subsets with 567 fitting, 36 selection, 35 calibration and 71 internal-test engines. Faults follow related-channel, temporally persistent mechanisms in raw sensor units. Conservative thresholds, fitting-only support bounds and a complete 41-condition audit limit healthy-sensor displacement, while detector recall remains low on drift, bias and noise. Predictor refinement uses active-path paired-prefix continuity, fitting-error tails and corruption-aware RUL supervision. A validation-selected candidate is frozen before interval calibration and test evaluation. On 707 official endpoints its capped RMSE is 11.08 cycles, MAE is 7.19, and severity-0.2 stuck-fault RMSE averages 12.11 over three corruption replicates. Internal clean RMSE is 12.07 over 484 prefixes from 71 engines. Observable interval scales are fitted independently of engine-level conformal calibration. Raw RMSE is 21.27; 150 engines with raw RUL above 125 account for 77.7\% of raw squared error. The monotone plateau stretch was available, but its non-increasing capped-RMSE guard rejected every positive value (alpha=0), so the raw-target bias is not reduced. A previously inspected severe prediction failure is reported explicitly. Under broadband sensor noise at 20\% of the within-regime scale, capped RMSE rises from \CurrentNoiseOfficialZero{} to \CurrentNoiseOfficialTwenty{} cycles with \CurrentNoiseCoverageTwenty{}\% empirical coverage; at 50\% it is \CurrentNoiseOfficialFifty{} cycles. Reused benchmarks, one training seed and simulated faults limit generalization claims.
